\emph{(1).} Since $k^{2} = 1$ and $I^{2} = -1$ commute, $(1 + I)^{-1} = \tfrac12(1 - I)$ and $(1 + Ik)(1 - I) = 1 + k - I(1 - k)$, so $V_{k} = p_{+} - Ip_{-}$; this is the element of \claimref{synthesis.wedge-reflections-compose-to-the-two-pi-rotor-so-geometric-modular-conjugations-give-interacting-spin-statistics}(iv), where $V_{k}^{2} = k$ is proved. Being a polynomial in $k$ with coefficients in $\mathrm{span}\{1, I\}$, it is defined whenever $k$ is, commutes with whatever commutes with $k$ and $I$, and is carried to itself by every map that fixes $k$ and $I$. For odd $a$, $p_{\pm}a = ap_{\mp}$, so $V_{k}aV_{k}^{-1} = (p_{+} - Ip_{-})a(p_{+} + Ip_{-}) = I(p_{+}ap_{-} - p_{-}ap_{+}) = Ia(p_{-} - p_{+}) = -Iak$. For two modes with generators $e_{1}, \dots, e_{4}$ and $k = e_{1}e_{2}e_{3}e_{4}$, $-Ie_{1}k = -Ie_{2}e_{3}e_{4}$, a product of three generators, not a combination of generators (checked in $\Cl(4,1)$, where the pseudoscalar is a central unit: the result has grade two, while $ce_{j}$ and $cIe_{j}$ have grades one and four). With more modes $k$ contains more factors. For one mode with generators $e_{1}, e_{2}$, $k$ is $I e_{1}e_{2}$ up to sign and $-Ie_{1}k = \pm e_{2}$, a generator: there the twist is the quarter turn of the mode.

\emph{(2).} $\omega^{2} = (-1)^{10\cdot 9/2} = -1$. Moving $\gamma_{p}$ through $\omega$ passes nine other generators, so $\omega\gamma_{p} = -\gamma_{p}\omega$, and $\omega$ commutes with products of two generators. From $N_{j} = v_{j}v_{j}^{*} = \tfrac12(1 - I\gamma_{2j-1}\gamma_{2j})$, $\gamma_{2j-1}\gamma_{2j} = -I(1 - 2N_{j})$, so $\omega = (-I)^{5}\prod_{j}(1 - 2N_{j}) = -I(-1)^{N}$. On $V_{+}$, which contains the empty state ($N = 0$; the projector $\prod_{j}v_{j}^{*}v_{j}$ is annihilated by every $v_{j}^{*}$, and $\omega$ times it is $-I$ times it), $\Gamma = 1$ and $\omega = -I$; on $V_{-}$, $\omega = +I$. Since $\omega\gamma_{p}\omega^{-1} = -\gamma_{p}$, conjugation by $\omega$ is the grading automorphism, of order two, while on $V$ the element itself has order four. On a state of mode number $N$, $\exp(-\tfrac{\pi}{2}I(5 - 2N)) = (-I)^{5}(-1)^{N} = -I(-1)^{N}$, so $\omega = \exp(-\tfrac{\pi}{2}IX)$ with $X = 5 - 2N$; with $n$ and $k$ the slot and pair occupations, $Y = (3n - 2k)/6$ and $B - L = 1 - 2k/3$, so $5(B - L) - 4Y = 5 - 2n - 2k = X$. The rotors of the ten are generated by the $45$ bivectors, $\omega$ commutes with all of them, and $\omega^{2} = -1$ is the other central element: $\omega$ generates the centre, of order four (classical). Since $(\gamma_{2j-1}\gamma_{2j})^{2} = -1$, $\exp(\tfrac{\pi}{2}\gamma_{2j-1}\gamma_{2j}) = \gamma_{2j-1}\gamma_{2j}$, and the two generators of one mode are both slot generators or both pair generators, so each of these bivectors is among the $21$ of Proposition \claimref{standard-model.automorphisms-preserve-slot-pair-split-are}. The statement about fifteen components: removing one component of $V_{+}$ does not change the action on the others.

\emph{(3).} $T$ is a product of five elements each linear in the generators, and the five pairs are disjoint, so $T$ is of grade five (\claimref{standard-model.majorana-element}); $\omega$ anticommutes with every element of odd grade by (2), so $\omega T\omega^{-1} = -T$. On $V_{+}$, $\omega = -I$ and $B$ is bilinear over $\mathrm{span}\{1, I\}$, so $B(\omega\psi, T\omega\varphi) = I^{2}B(\psi, T\varphi) = -B(\psi, T\varphi)$; the same computation applies to every bilinear form in two fields of the sixteen, and in the algebra every mass form has grade $1$, $3$ or $5$ and anticommutes with $\omega$. The element $z = \exp(\tfrac{\pi}{2}(\gamma_{1}\gamma_{2} - \gamma_{3}\gamma_{4})) = -\gamma_{1}\gamma_{2}\gamma_{3}\gamma_{4}$ is generated by a difference of slot mode bivectors, which commutes with $T$ and lies in $\su(2)_{L}$; it anticommutes with the four slot generators and commutes with the six pair generators and with $T$. So $\omega z$ commutes with the slot generators and $(\omega z)T(\omega z)^{-1} = -T$. On states $\exp(-2\pi IY) = (-1)^{n}\exp(\tfrac{2\pi}{3}Ik)$, which is $z$ times a central element of $\su(3)$. Every element of the group generated by the subalgebra fixing $T$ commutes with $T$; $\omega$ does not, so it is not in that group; that subalgebra is the Standard Model algebra, and a vev along $T$ breaks Pati-Salam to it, by clauses (2) and (3) of the cited claim. The bit of the neutral singlet is fixed at the stages that leave $\uu(5)$ (\claimref{synthesis.neutral-singlet-bit-decided-by-gauge-algebra-leaving-u5}), which is where $\omega$ is a gauge transformation.

\emph{(4).} The six multiplets of $V_{+}$ under the Standard Model algebra are pairwise inequivalent, so an element commuting with the algebra is a scalar of $\mathrm{span}\{1, I\}$ on each, and with square $-1$ the scalar is $\pm I$. The signed sums $\pm 6 \pm 3 \pm 3 \pm 2 \pm 1 \pm 1$ include $16$, $-16$ and $0$ (for instance $6 + 2 - 3 - 3 - 1 - 1$, the doublets against the singlets, which is $\omega z$ up to an overall sign). The sum of $6, 3, 3, 2, 1$ is $15$, odd, and changing a sign changes a signed sum by an even number, so every signed sum for fifteen components is odd; a sum of an odd number of odd integers is odd. $15n_{g} \equiv -n_{g} \pmod{16}$.
